% Standalone TikZ figure: XAS measurement principle
% Can be compiled separately or included via \input
\begin{figure}[htbp]
  \centering
  \begin{tikzpicture}[
    scale=0.9,
    atom/.style={circle, draw, fill=blue!20, minimum size=0.5cm,
                 inner sep=0pt, font=\tiny},
    absorber/.style={circle, draw, fill=red!40, minimum size=0.6cm,
                     inner sep=0pt, font=\tiny\bfseries},
    photon/.style={-{Stealth}, decorate,
                   decoration={snake, amplitude=1.5pt, segment length=6pt},
                   thick, red!70!black},
    electron/.style={-{Stealth}, thick, blue!70!black, dashed},
    scattered/.style={-{Stealth}, thick, green!50!black, dashed},
  ]
    % X-ray beam
    \draw[photon] (-3, 0) -- (-0.5, 0)
      node[midway, above=4pt, font=\small, text=red!70!black] {$h\nu$};

    % Absorber atom
    \node[absorber] (abs) at (0,0) {A};

    % Neighbor atoms
    \node[atom] (n1) at (2.2, 0.8) {N};
    \node[atom] (n2) at (2.0, -1.0) {N};
    \node[atom] (n3) at (-1.0, 2.0) {N};
    \node[atom] (n4) at (1.5, 2.2) {N};

    % Photoelectron wave (outgoing)
    \draw[electron] (abs.east) -- (n1.west)
      node[midway, above, font=\scriptsize] {$e^-$};
    \draw[electron] (abs.east) -- (n2.north west);

    % Scattered wave (returning)
    \draw[scattered] (n1.south west) -- (abs.north east);
    \draw[scattered] (n2.north) -- (abs.south east);

    % Interference annotation
    \draw[<->, thick, gray] (0, -2) -- (2.2, -2)
      node[midway, below, font=\small] {$R_j$};

    % Labels
    \node[font=\small, text=red!70!black, below left] at (-0.3, -0.5)
      {Absorber};
    \node[font=\small, text=blue!70!black, right] at (2.5, 0.8)
      {Nachbar};

    % Right side: schematic spectrum
    \begin{scope}[xshift=7cm, yshift=-0.5cm]
      \draw[->, thick] (0, 0) -- (4.5, 0)
        node[right, font=\small] {$E$};
      \draw[->, thick] (0, 0) -- (0, 3.5)
        node[above, font=\small] {$\mu(E)$};

      % Pre-edge
      \draw[thick, blue!60!black] (0.2, 0.5) -- (1.2, 0.6);

      % Edge jump
      \draw[thick, blue!60!black] (1.2, 0.6) -- (1.4, 2.5);

      % EXAFS oscillations
      \draw[thick, blue!60!black]
        (1.4, 2.5) .. controls (1.7, 2.8) and (1.9, 2.2) ..
        (2.1, 2.5) .. controls (2.3, 2.8) and (2.5, 2.3) ..
        (2.7, 2.5) .. controls (2.9, 2.7) and (3.1, 2.4) ..
        (3.3, 2.5) .. controls (3.5, 2.6) and (3.7, 2.45) ..
        (4.0, 2.5);

      % Annotations
      \draw[<->, gray, thick] (1.15, 0.6) -- (1.15, 2.5)
        node[midway, left, font=\scriptsize] {$\Delta\mu_0$};
      \node[font=\scriptsize] at (1.3, 0.2) {$E_0$};
      \draw[gray, dashed] (1.3, 0) -- (1.3, 0.6);

      \draw[decorate, decoration={brace, amplitude=4pt}]
        (1.5, 3.2) -- (4.0, 3.2)
        node[midway, above=4pt, font=\scriptsize] {EXAFS};

      \draw[decorate, decoration={brace, amplitude=4pt}]
        (1.1, 3.2) -- (1.5, 3.2)
        node[midway, above=4pt, font=\scriptsize] {XANES};
    \end{scope}
  \end{tikzpicture}
  \caption{Prinzip der Röntgenabsorptionsspektroskopie.
    Links: Das einfallende Photon erzeugt am Absorberatom (A) ein
    Photoelektron, das an den Nachbaratomen (N) gestreut wird.
    Die Interferenz zwischen ausgehender und rückgestreuter Welle
    moduliert den Absorptionskoeffizienten.
    Rechts: Schematisches Absorptionsspektrum $\mu(E)$ mit
    Kantensprung $\Delta\mu_0$, XANES- und EXAFS-Bereich.}
  \label{fig:xas-principle}
\end{figure}
